\documentclass[12pt, a4paper]{book}

% ─── Packages ───────────────────────────────────────────────────────────────
\usepackage[a4paper, margin=2.5cm]{geometry}
\usepackage{fontenc}
\usepackage{inputenc}
\usepackage{microtype}
\usepackage{graphicx}
\usepackage{xcolor}
\usepackage{titlesec}

\usepackage{enumitem}
\usepackage{booktabs}
\usepackage{array}
\usepackage{tabularx}
\usepackage{fancyhdr}
\usepackage{hyperref}
\usepackage{mdframed}
\usepackage{multicol}
\usepackage{parskip}
\usepackage[table]{xcolor}
\usepackage{textcomp}
\usepackage{titletoc}
% ─── Colours ────────────────────────────────────────────────────────────────
\definecolor{primary}{HTML}{2E7D32}      % deep green
\definecolor{accent}{HTML}{FF8F00}       % warm amber
\definecolor{lightgray}{HTML}{F5F5F5}
\definecolor{darkgray}{HTML}{424242}

% ─── Hyperref ───────────────────────────────────────────────────────────────
\hypersetup{
    colorlinks=true,
    linkcolor=primary,
    urlcolor=accent
}

% ─── Header / Footer ────────────────────────────────────────────────────────
\pagestyle{fancy}
\fancyhf{}
\fancyhead[LE,RO]{\small\textcolor{darkgray}{\leftmark}}
\fancyhead[RE,LO]{\small\textcolor{primary}{\textbf{Jorge's Recipe Book}}}
\fancyfoot[C]{\textcolor{darkgray}{\thepage}}
\renewcommand{\headrulewidth}{0.4pt}

% ─── Section / Chapter Styling ──────────────────────────────────────────────
\titleformat{\chapter}[display]
  {\normalfont\huge\bfseries\color{primary}}
  {}
  {0pt}
  {\Huge}
  [{\vspace{0.5ex}\titlerule[1.5pt]}]

\titleformat{\section}
  {\normalfont\Large\bfseries\color{primary}}
  {}
  {0pt}
  {}
  [{\vspace{0.2ex}\textcolor{accent}{\rule{\linewidth}{1pt}}}]

% ─── Custom Environments ─────────────────────────────────────────────────────

% Macro box (coloured info box)
\newmdenv[
  backgroundcolor=lightgray,
  linecolor=primary,
  linewidth=1.5pt,
  roundcorner=6pt,
  skipabove=10pt,
  skipbelow=10pt,
  innerleftmargin=12pt,
  innerrightmargin=12pt,
  innertopmargin=10pt,
  innerbottommargin=10pt
]{macrobox}

% Tip box
\newmdenv[
  backgroundcolor=accent!10,
  linecolor=accent,
  linewidth=1.5pt,
  roundcorner=6pt,
  skipabove=10pt,
  skipbelow=10pt,
  innerleftmargin=12pt,
  innerrightmargin=12pt,
  innertopmargin=10pt,
  innerbottommargin=10pt
]{tipbox}

% ─── Helper Commands ─────────────────────────────────────────────────────────
\newcommand{\recipetitle}[1]{%
  \section*{\textcolor{primary}{#1}}%
  \addcontentsline{toc}{section}{#1}%
}

\newcommand{\recipemeta}[3]{%
  \begin{macrobox}
    \begin{tabular}{lll}
      \textbf{\textcolor{primary}{Cost per serve:}} & #1 &
      \textbf{\textcolor{primary}{Calories:}} #2 \quad
      \textbf{\textcolor{primary}{Protein:}} #3
    \end{tabular}
  \end{macrobox}%
}

\newcommand{\ingredientheader}{\subsection*{\textcolor{accent}{Ingredients}}}
\newcommand{\methodheader}{\subsection*{\textcolor{accent}{Method}}}
\newcommand{\noteheader}{\subsection*{\textcolor{accent}{Notes}}}

% ════════════════════════════════════════════════════════════════════════════
\begin{document}

% ─── Title Page ──────────────────────────────────────────────────────────────
\begin{titlepage}
  \centering
  \vspace*{3cm}
  {\fontsize{48}{56}\selectfont\bfseries\textcolor{primary}{Jorge's}\par}
  \vspace{0.3cm}
  {\fontsize{32}{40}\selectfont\bfseries\textcolor{accent}{Recipe Book}\par}
  \vspace{0.5cm}
  {\large\textcolor{darkgray}{A 4-Week Meal Plan for a Healthier You}\par}
  \vspace{1.5cm}
  \textcolor{primary}{\rule{0.6\linewidth}{2pt}}
  \vspace{1.5cm}

  {\large\textcolor{darkgray}{
    12 recipes designed around your lifestyle ---\\[6pt]
    high protein, budget friendly, and easy to batch cook.
  }\par}

  \vfill
  {\small\textcolor{darkgray}{Christchurch, New Zealand \\ \the\year}}
\end{titlepage}

% ─── Table of Contents ───────────────────────────────────────────────────────
\tableofcontents
\clearpage

% ════════════════════════════════════════════════════════════════════════════
%  4-WEEK PLAN OVERVIEW
% ════════════════════════════════════════════════════════════════════════════
\chapter*{Your 4-Week Plan}
\addcontentsline{toc}{chapter}{Your 4-Week Plan}

\begin{tipbox}
  \textbf{How your week works:}\\
  \textbf{Monday \& Thursday} --- Flatmate cooks for you.\\
  \textbf{Tuesday \& Friday} --- You cook for both of you (6 serves each).\\
  \textbf{Wednesday \& Saturday} --- Leftovers from the night before.\\
  \textbf{Sunday} --- Youth group dinner. No cooking needed!
\end{tipbox}

\vspace{0.5cm}

\begin{tabularx}{\linewidth}{|l|X|X|}
  \hline
  \rowcolor{primary!20}
  \textbf{Week} & \textbf{Tuesday} & \textbf{Friday} \\
  \hline
  Week 1 & Slow Cooker Mexican Chicken & Beef Mince \& Lentil Shepherd's Pie \\
  \hline
  Week 2 & Turkey Mince Stir Fry & Slow Cooker Chicken \& Chickpea Stew \\
  \hline
  Week 3 & Beef \& Vegetable Chilli & Slow Cooker Pulled Pork Wraps \\
  \hline
  Week 4 & Lemon Herb Baked Salmon Tray & Turkey \& Lentil Bolognese \\
  \hline
\end{tabularx}

\vspace{1cm}

\subsection*{\textcolor{primary}{Breakfast Rotation}}

\begin{tabularx}{\linewidth}{|l|X|}
  \hline
  \rowcolor{primary!20}
  \textbf{Week} & \textbf{Breakfast Prep} \\
  \hline
  Week 1 & Egg Muffins \\
  \hline
  Week 2 & Overnight Oats \\
  \hline
  Week 3 & Breakfast Burritos \\
  \hline
  Week 4 & Frittata Slices \\
  \hline
\end{tabularx}

\vspace{1cm}

\begin{macrobox}
  \textbf{\textcolor{primary}{Daily Targets}} \\[4pt]
  \begin{tabular}{ll}
    Calories: & 2,000--2,200 kcal/day \\
    Protein:  & 150g/day \\
    Steps:    & 8,000/day \\
  \end{tabular}\\[6pt]
  \textit{Note: Two meals alone won't quite reach these targets. Add a protein shake, Greek yogurt, or a handful of nuts as a daily snack.}
\end{macrobox}

% ════════════════════════════════════════════════════════════════════════════
%  DINNER RECIPES
% ════════════════════════════════════════════════════════════════════════════
\chapter{Dinner Recipes}

% ─── Recipe 1 ────────────────────────────────────────────────────────────────
\recipetitle{1. Slow Cooker Mexican Chicken}
\recipemeta{\$2.20/serve}{470 kcal}{48g}

\ingredientheader
\begin{itemize}[leftmargin=1.5em]
  \item 1 kg chicken thighs
  \item 1 can black beans, drained and rinsed
  \item 1 can corn, drained
  \item 1 jar medium salsa (400g)
  \item 1 tsp cumin
  \item 1 tsp smoked paprika
  \item 1 tsp garlic powder
  \item \textbf{To serve:} rice, Greek yogurt, shredded cheese
\end{itemize}

\methodheader
\begin{enumerate}[leftmargin=1.5em]
  \item Place chicken thighs in the slow cooker.
  \item Add black beans, corn, salsa, and all spices. Stir to combine.
  \item Cook on \textbf{low for 6--8 hours} or high for 3--4 hours.
  \item Shred the chicken with two forks directly in the slow cooker and stir through the sauce.
  \item Serve over rice with a dollop of Greek yogurt and shredded cheese.
\end{enumerate}

\begin{tipbox}
  \textbf{Tip:} Makes 6 serves. Leftovers keep in the fridge for 3--4 days or freeze for up to 3 months.
\end{tipbox}

\clearpage

% ─── Recipe 2 ────────────────────────────────────────────────────────────────
\recipetitle{2. Beef Mince \& Lentil Shepherd's Pie}
\recipemeta{\$2.80/serve}{480 kcal}{42g}

\ingredientheader
\begin{itemize}[leftmargin=1.5em]
  \item 500g lean beef mince
  \item 1 cup red lentils, rinsed
  \item 2 cans diced tomatoes
  \item 1 onion, diced
  \item 2 carrots, diced
  \item 2 stalks celery, diced
  \item 1 tbsp Worcestershire sauce
  \item 1 cup beef stock
  \item 4 large kumara (sweet potato), peeled and cubed
  \item Salt, pepper, butter for mash
\end{itemize}

\methodheader
\begin{enumerate}[leftmargin=1.5em]
  \item Preheat oven to 180\textdegree C.
  \item Brown mince in a large oven-safe pan over medium-high heat. Drain excess fat.
  \item Add onion, carrots, and celery. Cook for 5 minutes until softened.
  \item Stir in lentils, diced tomatoes, Worcestershire sauce, and beef stock. Simmer for 20 minutes until lentils are soft.
  \item Meanwhile, boil kumara until tender, then mash with butter, salt, and pepper.
  \item Top the mince mixture with kumara mash. Bake for 25 minutes until golden.
\end{enumerate}

\begin{tipbox}
  \textbf{Tip:} The red lentils bulk out the mince significantly --- you'd never notice they're there, but they add fibre and extra protein cheaply.
\end{tipbox}

\clearpage

% ─── Recipe 3 ────────────────────────────────────────────────────────────────
\recipetitle{3. Turkey Mince Stir Fry}
\recipemeta{\$2.50/serve}{460 kcal}{44g}

\ingredientheader
\begin{itemize}[leftmargin=1.5em]
  \item 500g turkey mince
  \item 1 bag (500g) frozen stir fry vegetables
  \item 3 tbsp soy sauce
  \item 2 tbsp oyster sauce
  \item 1 tbsp sesame oil
  \item 3 cloves garlic, minced
  \item 1 tsp fresh ginger, grated
  \item \textbf{To serve:} jasmine rice or noodles
\end{itemize}

\methodheader
\begin{enumerate}[leftmargin=1.5em]
  \item Cook rice or noodles according to packet instructions.
  \item Heat a wok or large frying pan over high heat. Add a splash of oil.
  \item Cook turkey mince, breaking it up, until browned (about 5--7 minutes).
  \item Add garlic and ginger, stir fry for 1 minute.
  \item Add frozen vegetables and stir fry for 3--4 minutes.
  \item Pour over soy sauce, oyster sauce, and sesame oil. Toss to combine.
  \item Serve immediately over rice or noodles.
\end{enumerate}

\begin{tipbox}
  \textbf{Tip:} Works equally well with beef mince if turkey isn't available. Day-old rice works best for fried rice variants.
\end{tipbox}

\clearpage

% ─── Recipe 4 ────────────────────────────────────────────────────────────────
\recipetitle{4. Slow Cooker Chicken \& Chickpea Stew}
\recipemeta{\$2.50/serve}{500 kcal}{46g}

\ingredientheader
\begin{itemize}[leftmargin=1.5em]
  \item 1 kg chicken thighs
  \item 2 cans chickpeas, drained and rinsed
  \item 1 can diced tomatoes
  \item 1 can coconut milk
  \item 2 tbsp medium curry paste
  \item 2 cups baby spinach
  \item 3 cloves garlic, minced
  \item 1 tsp fresh ginger, grated
  \item \textbf{To serve:} rice or naan bread
\end{itemize}

\methodheader
\begin{enumerate}[leftmargin=1.5em]
  \item Add chicken thighs, chickpeas, diced tomatoes, coconut milk, curry paste, garlic, and ginger to the slow cooker.
  \item Stir to combine. Cook on \textbf{low for 7--8 hours}.
  \item In the last 10 minutes, stir through baby spinach until wilted.
  \item Serve over rice or with warm naan bread.
\end{enumerate}

\begin{tipbox}
  \textbf{Tip:} Use Panang or Tikka Masala curry paste for a medium heat. Rogan Josh will be spicier.
\end{tipbox}

\clearpage

% ─── Recipe 5 ────────────────────────────────────────────────────────────────
\recipetitle{5. Beef \& Vegetable Chilli}
\recipemeta{\$2.60/serve}{490 kcal}{45g}

\ingredientheader
\begin{itemize}[leftmargin=1.5em]
  \item 500g lean beef mince
  \item 1 can kidney beans, drained
  \item 1 can black beans, drained
  \item 2 cans diced tomatoes
  \item 1 can corn, drained
  \item 1 capsicum, diced
  \item 1 onion, diced
  \item 2 tsp chilli powder (medium)
  \item 1 tsp cumin
  \item 1 tsp smoked paprika
  \item 3 cloves garlic, minced
  \item \textbf{To serve:} rice or baked potato, Greek yogurt
\end{itemize}

\methodheader
\begin{enumerate}[leftmargin=1.5em]
  \item Brown mince in a large pot over medium-high heat. Drain excess fat.
  \item Add onion, capsicum, and garlic. Cook for 5 minutes.
  \item Add all remaining ingredients. Stir to combine.
  \item Simmer on low heat for 30--40 minutes, stirring occasionally.
  \item Serve over rice or a baked potato with Greek yogurt on top.
\end{enumerate}

\begin{tipbox}
  \textbf{Tip:} This freezes brilliantly --- make a double batch and freeze half in portions for a quick future dinner.
\end{tipbox}

\clearpage

% ─── Recipe 6 ────────────────────────────────────────────────────────────────
\recipetitle{6. Slow Cooker Pulled Pork Wraps}
\recipemeta{\$2.40/serve}{480 kcal}{50g}

\ingredientheader
\begin{itemize}[leftmargin=1.5em]
  \item 1 kg pork shoulder
  \item 1 tbsp smoked paprika
  \item 1 tsp cumin
  \item 1 tsp garlic powder
  \item 1 tsp brown sugar
  \item 1 cup chicken stock
  \item \textbf{To serve:} large flour tortillas, coleslaw, Greek yogurt
\end{itemize}

\methodheader
\begin{enumerate}[leftmargin=1.5em]
  \item Mix smoked paprika, cumin, garlic powder, and brown sugar together to form a dry rub.
  \item Rub the spice mixture all over the pork shoulder.
  \item Place in slow cooker and pour chicken stock around the base.
  \item Cook on \textbf{low for 8 hours}.
  \item Shred the pork using two forks. Stir through the cooking juices.
  \item Serve in warmed tortillas with coleslaw and a dollop of Greek yogurt.
\end{enumerate}

\begin{tipbox}
  \textbf{Tip:} Pork shoulder is one of the cheapest cuts available. Buy it on special and freeze --- it slow cooks perfectly from thawed.
\end{tipbox}

\clearpage

% ─── Recipe 7 ────────────────────────────────────────────────────────────────
\recipetitle{7. Lemon Herb Baked Salmon Tray}
\recipemeta{\$4.50/serve}{450 kcal}{42g}

\ingredientheader
\begin{itemize}[leftmargin=1.5em]
  \item 6 salmon fillets
  \item 1 bunch broccolini
  \item 1 bunch asparagus
  \item 250g cherry tomatoes
  \item 2 tbsp olive oil
  \item Zest and juice of 2 lemons
  \item 3 cloves garlic, minced
  \item Fresh dill or mixed herbs
  \item Salt and pepper
  \item \textbf{To serve:} roasted baby potatoes
\end{itemize}

\methodheader
\begin{enumerate}[leftmargin=1.5em]
  \item Preheat oven to 200\textdegree C.
  \item Toss broccolini, asparagus, and cherry tomatoes in olive oil, salt and pepper. Spread across a large baking tray.
  \item Place salmon fillets on top. Mix lemon zest, lemon juice, and garlic together and spoon over the salmon.
  \item Season with fresh dill, salt, and pepper.
  \item Bake for 20--25 minutes until salmon is cooked through and flakes easily.
  \item Serve with roasted baby potatoes.
\end{enumerate}

\begin{tipbox}
  \textbf{Tip:} Slightly pricier than other meals but worth it once a month. Watch for salmon on special at Pak'nSave or New World.
\end{tipbox}

\clearpage

% ─── Recipe 8 ────────────────────────────────────────────────────────────────
\recipetitle{8. Turkey \& Lentil Bolognese}
\recipemeta{\$2.30/serve}{460 kcal}{44g}

\ingredientheader
\begin{itemize}[leftmargin=1.5em]
  \item 500g turkey mince
  \item 1 cup red lentils, rinsed
  \item 2 cans diced tomatoes
  \item 2 tbsp tomato paste
  \item 1 onion, diced
  \item 2 carrots, diced
  \item 2 stalks celery, diced
  \item 3 cloves garlic, minced
  \item 1 tsp dried oregano
  \item 1 tsp dried basil
  \item 1 cup beef or chicken stock
  \item \textbf{To serve:} pasta or zucchini noodles, parmesan
\end{itemize}

\methodheader
\begin{enumerate}[leftmargin=1.5em]
  \item Brown turkey mince in a large pot. Drain excess fat.
  \item Add onion, carrots, celery, and garlic. Cook for 5 minutes until softened.
  \item Stir in tomato paste and cook for 1 minute.
  \item Add diced tomatoes, red lentils, stock, oregano, and basil.
  \item Simmer on low for 30 minutes until lentils are completely soft and sauce has thickened.
  \item Serve over pasta or zucchini noodles with parmesan.
\end{enumerate}

\begin{tipbox}
  \textbf{Tip:} For a lower carb option, zucchini noodles cut the calories significantly while keeping the meal filling.
\end{tipbox}

% ════════════════════════════════════════════════════════════════════════════
%  BREAKFAST RECIPES
% ════════════════════════════════════════════════════════════════════════════
\chapter{Breakfast Recipes}

% ─── Breakfast 1 ─────────────────────────────────────────────────────────────
\recipetitle{1. Egg Muffins}
\recipemeta{\$0.80/serve (2 muffins)}{320 kcal}{28g}

\ingredientheader
\begin{itemize}[leftmargin=1.5em]
  \item 10 eggs
  \item 100g diced bacon or ham
  \item 1 capsicum, finely diced
  \item Large handful of spinach, roughly chopped
  \item 100g cherry tomatoes, halved
  \item 50g feta or cheddar, crumbled or grated
  \item Salt, pepper, mixed dried herbs
\end{itemize}

\methodheader
\begin{enumerate}[leftmargin=1.5em]
  \item Preheat oven to 180\textdegree C. Grease a 12-hole muffin tin well.
  \item Whisk eggs together in a large bowl. Season with salt, pepper, and herbs.
  \item Stir in bacon, capsicum, spinach, tomatoes, and cheese.
  \item Pour mixture evenly into the muffin tin, filling each hole about three-quarters full.
  \item Bake for 18--22 minutes until set and lightly golden on top.
  \item Cool completely before storing in an airtight container in the fridge for up to 5 days.
\end{enumerate}

\begin{tipbox}
  \textbf{Reheat:} 60--90 seconds in the microwave. Eat 2 muffins per serving.
\end{tipbox}

\clearpage

% ─── Breakfast 2 ─────────────────────────────────────────────────────────────
\recipetitle{2. Overnight Oats}
\recipemeta{\$1.20/jar}{450 kcal}{35g}

\textit{Makes 5 jars --- prep all at once on Monday evening.}

\ingredientheader

\textit{Per jar:}
\begin{itemize}[leftmargin=1.5em]
  \item \textonehalf{} cup rolled oats
  \item \textonehalf{} cup Greek yogurt (full fat)
  \item 1 cup milk (or almond/oat milk)
  \item 1 scoop vanilla protein powder
  \item 1 tbsp chia seeds
  \item \textbf{Toppings (add in the morning):} banana, mixed berries, honey, granola
\end{itemize}

\methodheader
\begin{enumerate}[leftmargin=1.5em]
  \item Add oats, Greek yogurt, milk, protein powder, and chia seeds to a jar or container. Stir well.
  \item Seal and refrigerate overnight (or at least 4 hours).
  \item In the morning, give it a stir and add your toppings.
  \item Grab and go --- no heating required.
\end{enumerate}

\begin{tipbox}
  \textbf{Tip:} Keeps in the fridge for up to 5 days. Prep all 5 jars on Sunday or Monday to cover the whole week.
\end{tipbox}

\clearpage

% ─── Breakfast 3 ─────────────────────────────────────────────────────────────
\recipetitle{3. Breakfast Burritos}
\recipemeta{\$1.80/burrito}{480 kcal}{38g}

\textit{Makes 6 burritos --- freeze and microwave as needed.}

\ingredientheader
\begin{itemize}[leftmargin=1.5em]
  \item 8 eggs, whisked
  \item 200g chorizo or breakfast sausage, sliced
  \item 1 can black beans, drained and rinsed
  \item 1 capsicum, diced
  \item 1 onion, diced
  \item 6 large flour tortillas
  \item 80g cheddar cheese, grated
  \item Salsa and Greek yogurt to serve
  \item Salt, pepper, cumin, smoked paprika
\end{itemize}

\methodheader
\begin{enumerate}[leftmargin=1.5em]
  \item Cook chorizo in a large frying pan over medium heat until browned. Set aside.
  \item In the same pan, saut\'{e} capsicum and onion until softened (about 5 minutes).
  \item Add black beans and season with cumin, paprika, salt, and pepper. Cook for 2 minutes.
  \item Push vegetables to the side. Pour in whisked eggs and scramble until just cooked. Mix everything together.
  \item Lay out tortillas. Divide the filling evenly across all 6. Top with cheese.
  \item Roll tightly into burritos, folding in the sides as you go.
  \item Wrap each burrito individually in foil or cling wrap. Freeze.
\end{enumerate}

\begin{tipbox}
  \textbf{Reheat from frozen:} Remove foil, wrap in a damp paper towel, microwave for 2--3 minutes. Or reheat in foil in oven at 180\textdegree C for 15--20 minutes for a crispier result.
\end{tipbox}

\clearpage

% ─── Breakfast 4 ─────────────────────────────────────────────────────────────
\recipetitle{4. Frittata Slices}
\recipemeta{\$1.10/serve (2 slices)}{380 kcal}{32g}

\textit{Makes 8 slices --- keep in fridge, reheat as needed.}

\ingredientheader
\begin{itemize}[leftmargin=1.5em]
  \item 10 eggs
  \item 150g cooked chicken breast, diced
  \item 1 zucchini, grated (squeeze out excess moisture)
  \item 1 capsicum, diced
  \item Large handful of spinach
  \item 100g feta, crumbled
  \item 2 cloves garlic, minced
  \item 1 tsp dried Italian herbs
  \item Salt and pepper
  \item Olive oil for cooking
\end{itemize}

\methodheader
\begin{enumerate}[leftmargin=1.5em]
  \item Preheat oven to 180\textdegree C.
  \item Heat olive oil in a large oven-safe frying pan over medium heat. Saut\'{e} garlic, capsicum, and zucchini for 4--5 minutes. Add spinach and wilt.
  \item Whisk eggs in a large bowl with Italian herbs, salt, and pepper. Stir in chicken and feta.
  \item Pour egg mixture over the vegetables in the pan. Stir gently to distribute.
  \item Cook on the stovetop for 4--5 minutes until the edges begin to set.
  \item Transfer to the oven and bake for 15--20 minutes until fully set and golden on top.
  \item Cool, then slice into 8 portions. Store in an airtight container in the fridge for up to 5 days.
\end{enumerate}

\begin{tipbox}
  \textbf{Reheat:} 60--90 seconds in the microwave or a few minutes in a dry frying pan. Eat 2 slices per serving.
\end{tipbox}

% ════════════════════════════════════════════════════════════════════════════
%  NUTRITION REFERENCE
% ════════════════════════════════════════════════════════════════════════════
\chapter{Nutrition Reference}

\section*{Daily Targets}
\addcontentsline{toc}{section}{Daily Targets}

\begin{macrobox}
  \begin{tabular}{ll}
    \textbf{Daily Calories:} & 2,000--2,200 kcal \\
    \textbf{Daily Protein:}  & 150g \\
    \textbf{Daily Steps:}    & 8,000 \\
  \end{tabular}\\[8pt]
  Remember to log everything in \textbf{MyFitnessPal} and adjust your portions if your weight loss stalls or you feel overly hungry.
\end{macrobox}

\section*{Recipe Macros at a Glance}
\addcontentsline{toc}{section}{Recipe Macros at a Glance}

\begin{tabularx}{\linewidth}{|X|c|c|c|}
  \hline
  \rowcolor{primary!20}
  \textbf{Recipe} & \textbf{Cost/serve} & \textbf{Calories} & \textbf{Protein} \\
  \hline
  Mexican Chicken          & \$2.20 & 470 kcal & 48g \\
  \hline
  Shepherd's Pie           & \$2.80 & 480 kcal & 42g \\
  \hline
  Turkey Stir Fry          & \$2.50 & 460 kcal & 44g \\
  \hline
  Chicken \& Chickpea Stew & \$2.50 & 500 kcal & 46g \\
  \hline
  Beef \& Veg Chilli       & \$2.60 & 490 kcal & 45g \\
  \hline
  Pulled Pork Wraps        & \$2.40 & 480 kcal & 50g \\
  \hline
  Baked Salmon Tray        & \$4.50 & 450 kcal & 42g \\
  \hline
  Turkey \& Lentil Bolo    & \$2.30 & 460 kcal & 44g \\
  \hline
  \hline
  \rowcolor{accent!20}
  Egg Muffins (2)          & \$0.80 & 320 kcal & 28g \\
  \hline
  \rowcolor{accent!20}
  Overnight Oats           & \$1.20 & 450 kcal & 35g \\
  \hline
  \rowcolor{accent!20}
  Breakfast Burrito        & \$1.80 & 480 kcal & 38g \\
  \hline
  \rowcolor{accent!20}
  Frittata Slices (2)      & \$1.10 & 380 kcal & 32g \\
  \hline
\end{tabularx}

\vspace{1cm}

\begin{tipbox}
  \textbf{Hitting 150g protein daily on 2 meals:}\\
  Breakfast (28--38g) + Dinner (42--50g) = roughly 70--88g from meals alone. \\
  Bridge the gap with: 1 protein shake (\textasciitilde25g) + 200g Greek yogurt (\textasciitilde20g) + handful of almonds (\textasciitilde6g) = \textasciitilde51g extra. \\
  Total: \textbf{\textasciitilde140--140g} --- very close to target with minimal extra eating.
\end{tipbox}

% ════════════════════════════════════════════════════════════════════════════
%  SHOPPING TIPS
% ════════════════════════════════════════════════════════════════════════════
\chapter{Shopping Tips}
\addcontentsline{toc}{chapter}{Shopping Tips}

\section*{Budget Shopping in Christchurch}
\addcontentsline{toc}{section}{Budget Shopping in Christchurch}

\begin{itemize}[leftmargin=1.5em, itemsep=6pt]
  \item \textbf{Pak'nSave} for bulk staples --- rice, oats, canned goods, frozen vegetables, and eggs are significantly cheaper here.
  \item \textbf{Chicken thighs over breasts} --- roughly half the price per kilo, more flavour, and perfect for slow cooking.
  \item \textbf{Frozen vegetables} are just as nutritious as fresh and far cheaper. Keep a bag of frozen stir fry mix and frozen spinach always on hand.
  \item \textbf{Buy mince in bulk} when on special. Divide into 500g portions and freeze immediately.
  \item \textbf{Canned legumes} (chickpeas, kidney beans, black beans, lentils) are your best friend --- around \$1.50 a can and dramatically bulk out any meal.
  \item \textbf{Red lentils} are incredibly cheap per kilo and invisible in mince-based dishes --- they add fibre, protein, and volume.
  \item \textbf{Greek yogurt} --- buy the large tub (1kg) rather than individual portions. Cheaper and lasts the week.
  \item \textbf{Eggs} --- buy 18-packs at Pak'nSave for the best value.
\end{itemize}

\section*{Pantry Staples to Always Have}
\addcontentsline{toc}{section}{Pantry Staples to Always Have}

\begin{multicols}{2}
\begin{itemize}[leftmargin=1.5em, itemsep=4pt]
  \item Jasmine rice (2kg bag)
  \item Red lentils (1kg bag)
  \item Rolled oats (1kg)
  \item Canned diced tomatoes
  \item Canned chickpeas \& beans
  \item Chicken \& beef stock
  \item Soy sauce
  \item Oyster sauce
  \item Sesame oil
  \item Olive oil
  \item Smoked paprika
  \item Cumin
  \item Garlic powder
  \item Dried oregano \& basil
  \item Chilli powder
  \item Brown sugar
  \item Worcestershire sauce
  \item Tomato paste
  \item Large flour tortillas
  \item Protein powder (vanilla)
\end{itemize}
\end{multicols}

\vspace{1cm}
\begin{macrobox}
  \textbf{Weekly grocery target: \$60--\$100 NZD}\\[4pt]
  With pantry staples already stocked, your weekly fresh shop should sit comfortably in this range covering all breakfasts and your two cooking nights.
\end{macrobox}

\end{document}